\noindent\textbf{1.} A paragraph to put the equation below far off the screen: paragraph 1
says little, at some length, so that it takes a few lines at every width
the tests give the block, and the display stays out of view.

\noindent\textbf{2.} A paragraph to put the equation below far off the screen: paragraph 2
says little, at some length, so that it takes a few lines at every width
the tests give the block, and the display stays out of view.

\noindent\textbf{3.} A paragraph to put the equation below far off the screen: paragraph 3
says little, at some length, so that it takes a few lines at every width
the tests give the block, and the display stays out of view.

\noindent\textbf{4.} A paragraph to put the equation below far off the screen: paragraph 4
says little, at some length, so that it takes a few lines at every width
the tests give the block, and the display stays out of view.

\noindent\textbf{5.} A paragraph to put the equation below far off the screen: paragraph 5
says little, at some length, so that it takes a few lines at every width
the tests give the block, and the display stays out of view.

\noindent\textbf{6.} A paragraph to put the equation below far off the screen: paragraph 6
says little, at some length, so that it takes a few lines at every width
the tests give the block, and the display stays out of view.

\noindent\textbf{7.} A paragraph to put the equation below far off the screen: paragraph 7
says little, at some length, so that it takes a few lines at every width
the tests give the block, and the display stays out of view.

\noindent\textbf{8.} A paragraph to put the equation below far off the screen: paragraph 8
says little, at some length, so that it takes a few lines at every width
the tests give the block, and the display stays out of view.

\noindent\textbf{9.} A paragraph to put the equation below far off the screen: paragraph 9
says little, at some length, so that it takes a few lines at every width
the tests give the block, and the display stays out of view.

\noindent\textbf{10.} A paragraph to put the equation below far off the screen: paragraph 10
says little, at some length, so that it takes a few lines at every width
the tests give the block, and the display stays out of view.

\noindent\textbf{11.} A paragraph to put the equation below far off the screen: paragraph 11
says little, at some length, so that it takes a few lines at every width
the tests give the block, and the display stays out of view.

\noindent\textbf{12.} A paragraph to put the equation below far off the screen: paragraph 12
says little, at some length, so that it takes a few lines at every width
the tests give the block, and the display stays out of view.

\noindent\textbf{13.} A paragraph to put the equation below far off the screen: paragraph 13
says little, at some length, so that it takes a few lines at every width
the tests give the block, and the display stays out of view.

\noindent\textbf{14.} A paragraph to put the equation below far off the screen: paragraph 14
says little, at some length, so that it takes a few lines at every width
the tests give the block, and the display stays out of view.

\noindent\textbf{15.} A paragraph to put the equation below far off the screen: paragraph 15
says little, at some length, so that it takes a few lines at every width
the tests give the block, and the display stays out of view.

\noindent\textbf{16.} A paragraph to put the equation below far off the screen: paragraph 16
says little, at some length, so that it takes a few lines at every width
the tests give the block, and the display stays out of view.

\noindent\textbf{17.} A paragraph to put the equation below far off the screen: paragraph 17
says little, at some length, so that it takes a few lines at every width
the tests give the block, and the display stays out of view.

\noindent\textbf{18.} A paragraph to put the equation below far off the screen: paragraph 18
says little, at some length, so that it takes a few lines at every width
the tests give the block, and the display stays out of view.

\noindent\textbf{19.} A paragraph to put the equation below far off the screen: paragraph 19
says little, at some length, so that it takes a few lines at every width
the tests give the block, and the display stays out of view.

\noindent\textbf{20.} A paragraph to put the equation below far off the screen: paragraph 20
says little, at some length, so that it takes a few lines at every width
the tests give the block, and the display stays out of view.

\noindent\textbf{21.} A paragraph to put the equation below far off the screen: paragraph 21
says little, at some length, so that it takes a few lines at every width
the tests give the block, and the display stays out of view.

\noindent\textbf{22.} A paragraph to put the equation below far off the screen: paragraph 22
says little, at some length, so that it takes a few lines at every width
the tests give the block, and the display stays out of view.

\noindent\textbf{23.} A paragraph to put the equation below far off the screen: paragraph 23
says little, at some length, so that it takes a few lines at every width
the tests give the block, and the display stays out of view.

\noindent\textbf{24.} A paragraph to put the equation below far off the screen: paragraph 24
says little, at some length, so that it takes a few lines at every width
the tests give the block, and the display stays out of view.

The same numbered equation, far down the page:
\begin{equation}
  T_n = \frac{n(n+1)}{2}.
\end{equation}
and a line after it.
